\chapter*{前言}
\addcontentsline{toc}{chapter}{前言}
\markboth{前言}{前言}

市面上的面经大多是“题目 + 答案”的清单：你照着背，背完了，遇到稍微变形的题就又卡住。这本书想做一件不一样的事——\textbf{把“怎么想”讲清楚}。

Airbnb 的面试有一个鲜明的特点：它不喜欢考冷僻的算法套路，而喜欢考\textbf{贴近真实业务的场景题}，并且\textbf{会在你写完之后不断追加新需求}。这意味着，单纯记住某道题的最优解几乎没用——面试官真正想看的，是你能不能把问题建模清楚、能不能写出可以继续扩展的代码、能不能在系统设计里讲清楚每一个取舍背后的“为什么”。

所以本书的每一道题、每一个系统，都会围绕四种结构化的“盒子”展开：

\begin{itemize}
  \item \textbf{【设计思想】}（红框）：这道题/这个系统真正考的、可以迁移到别处的心法。
  \item \textbf{【陷阱】}（橙框）：最容易写错、最容易漏的边界。
  \item \textbf{【追问】}（蓝框）：面试官现场最常见的 follow-up，以及怎么接。
  \item \textbf{【要点】}（绿框）：复杂度与速记结论。
\end{itemize}

如果你只有两周时间，请直接翻到最后一章的冲刺计划，按表执行；如果你有更多时间，建议从头读——因为本书刻意把“设计思想”作为一条主线贯穿算法与系统设计两部分，读完你会发现：\emph{防止订单超卖的乐观锁}、\emph{中位数数据流的对顶堆}、\emph{搜索系统的分层缓存}，本质上都是同一批思维工具的不同投影。

祝面试顺利。

\vspace{1em}
\noindent\hfill —— SZY
